\subsection*{Independent reading of the discrete mixed-characteristic extension}

The complete 782-line text of
\texttt{arithmetic\_notes/B8\_discrete\_residue.tex}
was read on 18 November 2026.  Both B.META.20/B.146
and B.META.21/B.147 receive \textbf{PASS}.
The reader is GPT-6 Astra Pro in Pro mode; a separate
effort setting is not exposed.  This is a separate
same-model-family reading, with no priority claim.

\medskip\noindent\textbf{B.146: completion and Kummer theory.}
The exact Temkin and Stacks completion statements
already opened for B.139 apply to the present DVR.
Neither requires a finite residue field.
Relative algebraic closedness and finite Galois
linear disjointness preserve extension degrees;
the unique extended henselian valuation gives
full faithfulness.  The finite algebraic part
therefore supplies both a basis and dense scalar
points in every finite completed extension.

For the finitely generated residue field, the
intersection of all $p^n$-power subfields is its
finite constant field.  A divisor on a normal
proper model proves this: all coefficients are
divisible by every $p^n$, hence vanish, and
normality and properness make the function
constant.  No perfectness is assumed.

For $m>e/(p-1)$, logarithm and exponential are
inverse on the indicated deep unit group.
Repeated $p$-th powers carry $U_1$ into it
uniformly.  If $b^{p^n}$ is a deep unit, valuation
and reduction force $b\in U_1$, also when the
residue field is imperfect.  This proves the
two inclusions in (META.DV4) and the full induced
$p$-power topology.  Completeness gives the
injection into the inverse-limit Kummer group.
The possible $\varprojlim^1$ term vanishes for
the finite roots-of-unity groups.  Bounded
continuous cocycles and coboundaries justify
tensoring to rational coefficients.

The resulting map $\iota_M$ is injective,
extension-compatible and equivariant by its
explicit exponential formula.  Surjectivity
onto all cohomology is never asserted.
A zero rational Kummer class has value zero
and residue in the finite constant field.
After removing a prime-to-$p$ root of unity
and taking a bounded $p$-power it lies in
the deep-unit group, where injectivity makes
it one.  The converse is immediate for torsion.
The cyclotomic valuation denominators bound
the $p$-power roots of unity, including at two;
the other roots inject into the finite constant
field.  This proves exactly the stated finite
root-of-unity kernel.

For $T=M^{\mathrm{nr}}\cap\overline{\mathbb Q}_p$,
unramified constant extensions remain etale.
Every finite constant subextension has ramification
index dividing the fixed $e$; after maximal
unramified base change its degree equals that
index by an Eisenstein uniformizer polynomial.
The union has bounded finite subextension degrees,
so $T/\mathbb Q_p^{\mathrm{nr}}$ is finite.
Compactness makes the inertia image closed,
and Galois correspondence identifies it with
the inertia group of a finite local field.
Its nontrivial cyclotomic action gives the
restriction injection.  Its strictly decreasing
derived series transfers by surjectivity.
The proof uses known wild freeness only for
that ordinary finite local quotient, never
for the original imperfect-residue field.

\medskip\noindent\textbf{B.147: arithmetic inputs.}
The complete primary Stacks Tag 0BUA proof
was independently reopened.  It applies to an
arbitrary DVR.  Adjoining prime-to-$p$ roots
of a uniformizer over the strict henselization
gives the full tame quotient in this application.
The residue field is correctly retained as
$k^{\mathrm{sep}}$, which can be imperfect.

For derived fields, the two-step cohomology
filtration has trivial conjugation on both
pieces.  Thus $(g-1)^2=0$; the finite orbit
of the specific Kummer vector forces it to
be fixed in characteristic zero.  One finite
root-of-unity group gives the uniform exponent.
Strict derived inclusions supply the nonzero
class via $b$ and $1-b$.  The two complex
grids then recover a monic rational function
by the bounded-degree absolute-square identity
and its nonzero additive period.

For the nilpotent union, the tame pro-$\ell$
factor is central.  Choosing
$\ell\nmid ep(p-1)$ gives a moving element
with trivial cyclotomic action.  The radical
group is torsion-free pro-$p$, whereas the
image of the lower-central intersection is
finite.  Their intersection kills the required
commutators.  The explicit cocycle calculation
gives injectivity through the cyclotomic quotient.
This yields exactly the stable-kernel implication
needed for the same nonzero-class and two-grid
arguments, at every prime.

\medskip\noindent\textbf{B.147: the local analytic step.}
The good projection fibre is split over the
local field of definition before the cyclic
extension is chosen.  No later splitting
extension can reduce its degree $q>d+1$.
Taking the root that makes the integral
Kummer image surjective occurs inside the
kernel's fixed field; the required finite
local enlargement preserves its normality.
B.146 supplies an algebraic basis of $L$
and dense algebraic scalar parameter points.

The class $\theta_t$ inflates from
$G_L/(N_1\cap G_L)$.  Normality of $N_1$
and of $G_L$ makes conjugation by $N_1$
trivial on this quotient, and its Tate
coefficient action is also trivial.
On a small actual ball all ratios are deep
units, so subtracting errors gives exactly
(META.DV11).  The injective equivariant
$\iota_L$ makes the scalar invariant under
$N_1$, whose image is all of
$\operatorname{Gal}(L/M_1)$; it belongs to $M_1$.
Applying $\tau^i$ and subtracting the original
identity gives (META.DV12) with its written
sign.  The centres and boundary coefficients
belong to $M_1$.  Exponentiation is used only
after shrinking the near-one ratios into its
convergence ball.  No logarithm on arbitrary
residue-unit classes is presumed.

Minimizing the sum of gcd degrees and preserving
nonzero subresultants makes each reduced degree
locally constant on actual points.  The stated
$2e_i$-value interpolation is invertible by
the degree bound, coprimality and monicity.
It constructs analytic reduced ratios without
assuming the actual polynomial analytic.
The scalar blocks, root types and divisor
cocycle then have the precise B.136 inputs.
A constant root is algebraic after an actual
specialization, and its persistence violates
the dimension bound for a proper actual fibre.
A negative type-zero multiplicity would give
$q-1>d$ distinct denominator roots, impossible
after a degree-preserving specialization.

Only finitely many analytic minors and
subresultants must be avoided.  Their
complements are open and dense by the
characteristic-zero power-series identity
principle; their finite intersection meets
the dense algebraic scalar set.  Neither
local compactness nor Baire's theorem enters.
The analytic denominator gives the effective
divisor directly.  B.135, finite rational
interpolation, integral residues and the
normalized two-grid test recover the full
polynomial and remove its common factor.
The graph-degree proof is algebraic and
transfers with no new local-field hypothesis.
Continuity removes poles and exceptional
point values as in the previously checked
B.136 and B.140 proofs.

\medskip\noindent\textbf{Descent and exact additional scope.}
Finite-presentation descent of etale objects,
morphisms and finite covering families justifies
the compatible finite-slice limit over an
infinite algebraic extension.  The actual
point map retains equivariance and descends
the unique reconstructed scheme morphism.
Finite composita handle all finite extensions,
including those of the first two derived fields.

Kuhlmann Theorem 9.44 and its full proof
were independently opened.  In the Gauss
examples both candidate henselian valuations
are nontrivial of rank one.  Inequivalence
would imply independence and contradict that
theorem on a non-separably-closed field.
Equivalence would identify their residue
transcendence degrees, contradicting $s>0$
versus zero.  Thus the extra family is
strict even without assuming in advance
that a field isomorphism preserves valuations.
The proof correctly stops short of the
all-proper-normal-subgroup assertion of B.143.

\medskip\noindent\textbf{Primary sources opened for this event.}
Stacks Project, Tag 0BUA, Lemma 58.13.5,
entire statement and proof:
\url{https://stacks.math.columbia.edu/tag/0BUA}.
Kuhlmann, \emph{Valuation theory}, Chapter 9,
Theorem 9.44, printed page 240, statement
and proof:
\url{https://www.fvkuhlmann.de/bookch9.pdf}.
The exact Temkin, Stacks 09EJ, JTWY and CHW
inputs were previously opened and recorded
in the B.139--140 and B.131--136 audits;
their hypotheses have been checked in this
new application.  No complete re-proof of
those external inputs is claimed here.
